% Table: Robustness of earnings estimates (Chapter 4 appendix)
\begin{table}[htbp]
\centering
\caption{Robustness of IV earnings estimates to alternative specifications}
\label{tab:robustness}
\begin{tabular}{l ccc}
\toprule
                                          & First stage  & IV: log      & First-stage \\
                                          & ($\pi_1$)    & earnings     & $F$         \\
\midrule
1. Baseline (linear controls)             & 0.057\sym{***}  & 0.135\sym{*}   & 843 \\
                                          & (0.002)         & (0.069)         &     \\
\addlinespace
2. Saturated, region as London--SE indicator (141 cells) & 0.058\sym{***}  & 0.152\sym{**}  & 891 \\
                                          & (0.002)         & (0.072)         &     \\
\addlinespace
3. School-specific linear trends          & 0.049\sym{***}  & 0.125           & 804 \\
                                          & (0.002)         & (0.097)         &     \\
\addlinespace
4. Region $\times$ year FE                & 0.058\sym{***}  & 0.120\sym{*}   & 760 \\
                                          & (0.002)         & (0.071)         &     \\
\addlinespace
5. Saturated with KS2 attainment          & 0.057\sym{***}  & 0.145\sym{**}  & 823 \\
                                          & (0.002)         & (0.068)         &     \\
\addlinespace
6. Fine saturated (745 cells)             & 0.058\sym{***}  & 0.159\sym{**}  & 890 \\
                                          & (0.002)         & (0.074)         &     \\
\addlinespace
7. DDML                                   & ---             & 0.128\sym{*}   & --- \\
                                          &                 & (0.068)         &     \\
\addlinespace
School FE                                 & \multicolumn{3}{c}{All specifications} \\
Cohort FE                                 & \multicolumn{3}{c}{All specifications} \\
Observations                              & \multicolumn{3}{c}{368,260} \\
\bottomrule
\end{tabular}
\begin{minipage}{\textwidth}
\smallskip
\footnotesize
\textit{Notes:} Each row reports the first-stage coefficient, the IV estimate of the effect of A-level economics on log earnings at age 28, and the first-stage $F$-statistic. Row~1 is the baseline with linear individual controls. Row~2 uses the saturated covariate specification with region entered as a binary London--SE indicator (sex $\times$ ethnicity group $\times$ FSM $\times$ GCSE tercile $\times$ London-SE); the main-text saturated specification instead uses the three-category region (London/South East, rest of England, missing). Row~3 adds school-specific linear time trends. Row~4 replaces cohort FE with region $\times$ year FE. Row~5 replaces GCSE attainment with KS2 (age 11) attainment in the saturated cells. Row~6 uses the fine saturated specification, with GCSE vigintiles in place of terciles. Row~7 uses double/debiased machine learning. Standard errors clustered at the school level in parentheses. Sample: students in switcher schools, GCSE cohorts 2002--2008. The observation count refers to the baseline; the saturated specifications use slightly fewer observations because small cells drop (cf.\ Table~\ref{tab:main} column~4).

\smallskip
\noindent
All first-stage $F$-statistics exceed 104.7, the threshold above which the conventional critical value of 1.96 remains valid under the $tF$ procedure of \citet{lee2022valid}. No adjustment is therefore required in any specification. The $F$-statistics are cluster-robust throughout: with a single instrument, each equals the square of the instrument's cluster-robust $t$-statistic.

\smallskip
\noindent
Row~7 reports no first-stage coefficient or $F$-statistic because DDML does not estimate a first stage in the two-stage least squares sense. Following \citet{chernozhukov2018dml}, the conditional expectations of the treatment and the instrument given the covariates are estimated by lasso with cross-fitting, and the treatment effect is recovered from the resulting orthogonalised moment condition. There is accordingly no single coefficient on availability to report, and conventional weak-instrument diagnostics do not apply. For reference, the linear first stage estimated on the same sample is reported in row~1.

\smallskip
\noindent
\sym{*} $p<0.10$, \sym{**} $p<0.05$, \sym{***} $p<0.01$.
\end{minipage}
\end{table}
